\documentclass[11pt,a4paper,oneside,openany]{book}

\usepackage[utf8]{inputenc}
\usepackage[T1]{fontenc}
\usepackage[english]{babel}
\usepackage{csquotes}
\usepackage[charter]{mathdesign}
\usepackage{amsmath}
\usepackage{lipsum}
\usepackage{amsthm}
\usepackage{mathtools}
\usepackage{physics}
\usepackage{graphicx}
\usepackage{enumitem}
\usepackage{tikz-cd}
\usepackage[margin=2cm]{geometry}
\usepackage[backend=biber,style=numeric,sorting=none]{biblatex}
\addbibresource{bib.bib}
\usepackage[colorlinks=true,linkcolor=black,citecolor=black,urlcolor=black]{hyperref}

\newtheorem{theorem}{Theorem}[chapter]
\newtheorem{definition}{Definition}[section]
\newtheorem{lemma}[definition]{Lemma}
\newtheorem{proposition}[definition]{Proposition}
\newtheorem{corollary}[definition]{Corollary}
\theoremstyle{remark}
\newtheorem{remark}[definition]{Remark}
\newtheorem{example}[definition]{Example}

\let\mathcal\relax
\newcommand{\mathcal}[1]{\text{\usefont{OMS}{cmsy}{m}{n}#1}}
\DeclareMathOperator{\Spec}{Spec}
\DeclareMathOperator{\Ad}{Ad}
\DeclareMathOperator{\ad}{ad}
\newcommand{\poisson}[2]{\{#1,#2\}}
% ============================================================
% Comandi Matematici
% ============================================================

% --- Insiemi numerici ---
\newcommand{\R}{\mathbb{R}}
\newcommand{\C}{\mathbb{C}}
\newcommand{\N}{\mathbb{N}}
\newcommand{\K}{\mathbb{K}}
\newcommand{\A}{\mathfrak{A}}
\newcommand{\B}{\mathcal{B}}

% --- Spazio di Hilbert e spazio proiettivo ---
% ATTENZIONE: \H è già un comando di LaTeX (accento ungherese, es. \H{o}=ő).
% Il renewcommand lo sovrascrive: nella tesi non ti servirà mai quell'accento,
% quindi è sicuro, ma se preferisci non toccare comandi di sistema usa \Hil
% al posto di \H ovunque sotto.
\renewcommand{\H}{\mathcal{H}}
\newcommand{\PH}{\mathbb{P}\H}                    % P H, spazio proiettivo

% --- Algebre C* ---
\newcommand{\Cstar}{C^*}                          % da usare in math mode: $\Cstar$
\newcommand{\Csalg}{$C^*$-algebra}                % da usare nel testo: "è una \Csalg\ commutativa"
\newcommand{\Csalgs}{$C^*$-algebre}               % plurale
\newcommand{\BH}{B(\H)}                           % operatori limitati su H
\newcommand{\Cz}[1]{C_0(#1)}                      % C_0(M)
\newcommand{\Cinf}[1]{C^\infty(#1)}               % C^infty(M)
\newcommand{\Cinfc}[1]{C^\infty_c(#1)}            % C^infty_c(M)
\newcommand{\Mat}[2]{M_{#1}(#2)}                  % M_n(C), es. \Mat{2}{\C}

% --- Algebre di Lie (fraktur generico) ---
\newcommand{\mf}[1]{\mathfrak{#1}}                % \mf{g}, \mf{su}(2), \mf{u}(n), \mf{X}(M)

\newcommand{\qcomm}[2]{\dfrac{1}{i\hbar}[#1,#2]}  % (1/ih)[A,B]  -- condizione di Dirac, Qubit IV

% --- Notazione bra-ket ---
\newcommand{\melem}[3]{\langle#1|#2|#3\rangle}    % elemento di matrice <psi|A|psi>
\newcommand{\expect}[2]{\langle#1\rangle_{#2}}    % <A>_psi

% --- Campo vettoriale hamiltoniano ---
\newcommand{\Ham}[1]{X_{#1}}                      % \Ham{f}, \Ham{f_H}, \Ham{f_A}

% --- Simboli e operatori vari ---
\newcommand{\Id}{\mathbb{1}}
\renewcommand{\Re}{\mathrm{Re}}                   % la tesi usa Re/Im in tondo, non ℜ/ℑ
\renewcommand{\Im}{\mathrm{Im}}
\newcommand{\LC}{\varepsilon_{ijk}}               % simbolo di Levi-Civita (indici fissi ijk)
\DeclareMathOperator{\Ran}{Ran}
\DeclareMathOperator{\Ker}{Ker}
\DeclareMathOperator{\Span}{span}
\DeclareMathOperator{\sing}{sing}
\setcounter{chapter}{1}
\begin{document}

\noindent This chapter establishes the geometric and algebraic language of classical mechanics in an invariant form. Moving beyond the traditional Newtonian approach, we adopt a coordinate-free framework on smooth manifolds. The logical sequence of the chapter is as follows: first, we introduce Lie algebras (relying on Humphreys \cite{humphreysIntroductionLieAlgebras2010} and Lee \cite{leeIntroductionSmoothManifolds2013}) to provide the formal infrastructure for the algebraic brackets; second, using the foundations of symplectic geometry (following Abraham and Marsden \cite{abrahamFoundationsMechanics2008}), we identify $C^\infty(M)$ as the natural space for observables. In this context, Poisson brackets endow the space with a Lie algebra structure. Finally, by applying the Gelfand theorem, we recognize $C_0(M)$ as a commutative $C^*$-algebra. This leads to our primary structural result, classical mechanics formally coincides with the theory of commutative $C^*$-algebras.

\section{A Brief Introduction to Differential Geometry}

The geometric formulation of classical mechanics replaces the traditional Euclidean space with a coordinate-free framework on smooth manifolds. This abstraction is physically justified by holonomic constraints, which curve the configuration space and remove any global vector space structure. This section introduces the topological and differential tools required to extend ordinary calculus to such curved spaces.

\begin{definition}[Topological manifold]\label{def:top-man}
	A topological manifold of dimension $n$ is a topological space $M$ that is Hausdorff, second countable, and locally Euclidean of dimension $n$: every point $p \in M$ has an open neighborhood homeomorphic to an open subset of $\R^n$. We write $\dim M = n$.
\end{definition}

Physical configuration spaces provide the most natural examples of such topological spaces. For instance, a simple pendulum constrained to move in a plane has the unit circle $S^1$ as its configuration space, while the configurations of a free rigid body are described by the special orthogonal group $SO(3)$. Neither of these spaces is globally homeomorphic to a Euclidean space. To perform calculus, we must cover the manifold with open sets that locally resemble flat space.

\begin{definition}[Chart, atlas]\label{def:chart-atlas}
	A chart on $M$ is a pair $(U,\varphi)$, with $U \subseteq M$ open and $\varphi: U \to \varphi(U) \subseteq \R^n$ a homeomorphism. Two charts $(U,\varphi)$, $(V,\psi)$ are $C^\infty$-compatible if the transition maps
	\[
	\psi \circ \varphi^{-1} : \varphi(U \cap V) \to \psi(U \cap V)
	\]
	is a $C^\infty$-diffeomorphism whenever $U \cap V \neq \emptyset$. An atlas is a collection of pairwise compatible charts covering $M$. The union of all atlases compatible with a particular one on M is called the maximal atlas.
\end{definition}

\begin{remark}
	An equivalent formulation of Definition \ref{def:chart-atlas} could have been state using $\varphi \circ \psi^{-1}$ that is a diffeomorphism being the inverse of $\psi \circ \varphi^{-1}$. 
\end{remark}

The compatibility condition ensures that if two observers describe the same physical region using different local coordinates, the transformation between their descriptions is smooth. This allows for a globally consistent notion of differentiability.

\begin{definition}[Smooth manifold]\label{def:smooth-man}
	A smooth manifold is a pair $M \equiv (M,\mathcal{A})$, where $M$ is a n-dimensional topological manifold and $\mathcal{A}$ is a maximal atlas on M referred to as a smooth structure. 
\end{definition}

The configuration space is a smooth manifold. Observables are scalar functions on it, while symmetries and coordinate transformations are mappings between manifolds. Since these spaces generally lack a global Euclidean structure, smoothness must be checked locally, pulling the map back through charts into Euclidean space, where ordinary multivariable calculus applies.

\begin{definition}[Smooth map, diffeomorphism]\label{def:diffeo}
	Let $M$, $N$ be smooth manifolds, respectively of dimension $m, n \geq 1$. A map $f: M \to N$ is smooth at $p \in M$ if there exist charts $(U,\varphi)$ around $p$ and $(V,\psi)$ around $f(p)$, with $f(U) \subseteq V$, such that
	\[
	\psi \circ f \circ \varphi^{-1} : \varphi(U) \to \psi(V)
	\]
	is $C^\infty$ in the topology of $\R^m$. A function $f$ is smooth if it is such at every point, while it is a diffeomorphism if it is a smooth bijection with smooth inverse. We write $C^\infty(M) := C^\infty(M;\R)$.
\end{definition}
\medskip

The concept of manifold itself does not carry a vector space structure. Despite that, we can always fix a point and define a tangent vector as a derivation. This creates a vector space at every point of M called \textit{tangent space}.

\begin{definition}[Tangent vector]\label{def:tan-vector}
	Let $M$ be a smooth manifold as in Definition \ref{def:smooth-man} and let $p \in M$. A tangent vector at $p$ is a linear map $v : C^\infty(M) \to \R$ satisfying the Leibniz rule
	\[
	v(fg) = f(p)\,v(g) + g(p)\,v(f), \qquad f,g \in C^\infty(M).
	\]
	The set of tangent vectors at $p$ is the tangent space at a point $p$ denoted as $T_pM$.
\end{definition}

\begin{remark}\label{remark:vec-der}
	Equivalently, a smooth curve $\gamma:(-\varepsilon,\varepsilon)\to M$ with $\varepsilon > 0$ and $\gamma(0)=p$ induces the derivation $v_\gamma(f) := \frac{d}{dt}\Big|_{t=0} f(\gamma(t))$, as per definition \ref{def:tan-vector} corresponds to one of such maps. It can be proved that two curves induce the same tangent vector if and only if their local representatives agree to first order at $p$. This equivalent identification $T_pM \cong \{\text{curves through } p\}/\!\sim$ is used at a later stage.
\end{remark}

Most of the times it is convenient to write statements locally on a specific chart. This permits one to use the rules of calculus on manifolds in order to simplify computations.

\begin{definition}[Coordinate basis]\label{def:coord}
	Given a chart $(U,\varphi)$ around $p \in M$ n-dimensional smooth manifold, with $\varphi : U \to V \subseteq \R^n$ a homeomorphism. Let $r^i : \R^n \to \R$ with $i=1,...,n$ the standard coordinates of $\R^n$ which represent the canonical projections and are linearly independent.
	Define the coordinate basis as:
	\[
		x^i := r^i \circ \varphi : U \to \R \qquad i = 1,...,n
	\]
	the vectors
	\[
	\frac{\partial}{\partial x^i}\bigg|_p \in T_pM, \qquad
	\frac{\partial}{\partial x^i}\bigg|_p(f) := \frac{\partial (f\circ\varphi^{-1})}{\partial r^i}\bigg|_{\varphi(p)}
	\]
\end{definition}

\begin{proposition}
	A coordinate basis formed as in Definition \ref{def:coord}, forms a basis of $T_pM$. Hence $\dim T_pM = n$, and $v = v^i\, \partial_i |_p$ using the Einstein summation convention.
\end{proposition}

We can now assemble the tangent spaces at every point of $M$ into a single geometric object, endowed with its own smooth structure. In literature, it is used as the natural stage for the Lagrangian description of a mechanical system.

\begin{definition}[Tangent space]\label{def:tan-space}
	Let $M$ be a smooth manifold as in Definition \ref{def:smooth-man}, the tangent space of $M$ is
	\[
	TM := \bigsqcup_{p \in M} T_pM,
	\]
	it is endowed with the canonical projection $\pi: TM \to M$, $\pi(v) = p$ for $v \in T_pM$. 
	
	%A coordinate basis $(U,\varphi=(x^1,\dots,x^n))$ on $M$ induces a counterpart on $\pi^{-1}(U)$ via $v = v^i \partial_i|_p \mapsto (x^1(p),\dots,x^n(p),v^1,\dots,v^n)$, giving $TM$ a smooth structure of dimension $2n$.
\end{definition}

A single tangent vector specifies a generalized velocity at a specific configuration. To model global dynamical systems, we must assign a velocity vector to every point of the manifold in a smooth manner, which leads to the notion of a vector field.

\begin{definition}[Vector field]\label{def:vector-field}
	Let $M$ be a smooth manifold as in Definition \ref{def:smooth-man}, a vector field on $M$ is a smooth map $X: M \to TM$ with $\pi \circ X = \mathrm{id}_M$. We denote by $\mathfrak{X}(M)$ the $C^\infty(M)$-module of vector fields on $M$.
\end{definition}

The tangent space captures the kinematics of a system: it is the space of generalized velocities. The Hamiltonian framework also requires generalized momenta, which are dual to velocities and are modeled as linear functionals on the tangent space.

\begin{definition}[Covector]\label{def:covector}
	Let $M$ be a smooth manifold, the cotangent space at $p \in M$ is the dual vector space $T_p^*M := (T_pM)^*$. For $f \in C^\infty(M)$, its differential at $p$ is the covector
	\[
	df_p \in T_p^*M, \qquad df_p(v) := v(f), \quad v \in T_pM.
	\]
\end{definition}

\begin{remark}
	In the coordinate basis, $\{dx^1|_p,\dots,dx^n|_p\}$ is dual to $\{\partial/\partial x^i|_p\}$, and $df_p = \frac{\partial f}{\partial x^i}\Big|_p\, dx^i|_p$.
\end{remark}

The same construction that produced the tangent space from the individual tangent spaces can be repeated for their duals.

\begin{definition}[Cotangent space]\label{def:cot-space}
	Let $M$ be a smooth manifold, the cotangent space is
	\[
	T^*M := \bigsqcup_{p \in M} T_p^*M,
	\]
	it is endowed with the canonical projection $\pi: T^*M \to M$.
	%the induced smooth structure of dimension $2n$. A smooth section $\alpha: M \to T^*M$ is a differential $1$-form. The space of $1$-forms is $\Omega^1(M)$.
\end{definition}

The next proposition follows from the Vector Bundle Construction Lemma \cite[Lemma 10.6]{leeIntroductionSmoothManifolds2013} that will not be proved here. In summary it states that given a surjective local homeomorphism $\pi : X \to Y$ from a topological space $X$ and a smooth manifold $Y$, there exists a unique smooth structure on X, such that $\pi$ becomes a local diffeomorphism. 

\begin{proposition}[$TM$ and $T^*M$ as smooth manifolds]\label{prop:smooth-tm}
	Let $M$ be a smooth manifold of dimension $n$. Then $TM$ and $T^*M$, endowed
	with the natural projections $\pi_{TM} : TM \to M$ and $\pi_{T^*M} : T^*M \to M$,
	are smooth manifolds of dimension $2n$.
\end{proposition}

It is clear from Proposition \ref{prop:smooth-tm} that for $TM$ a coordinate basis $(U,\varphi=(x^1,\dots,x^n))$ on $M$ induces a counterpart on $\pi^{-1}(U)$ via $v = v^i \partial_i|_p \mapsto (x^1(p),\dots,x^n(p),v^1,\dots,v^n)$. The same construction can be extended to $T^*M$.

\noindent Now, we generalize the pointwise notion of covector to multilinear objects of arbitrary order, the building blocks from which differential forms are constructed.

\begin{definition}[Differential $k$-covectors]\label{def:k-covector}
	Let $M$ be a smooth manifold of dimension $n$ and let $k \geq 0$. For each $p \in M$, we denote by $\Lambda^k(T_p^*M)$ the vector space of alternating $k$-covectors at $p$ that are $k$-linear alternating maps on $T_pM$. The disjoint union
	\[
	\Lambda^k(T^*M) := \bigsqcup_{p \in M} \Lambda^k(T_p^*M),
	\]
	is the \emph{space of alternating $k$-covectors}.
\end{definition}

Using the same construction seen above in Proposition \ref{prop:smooth-tm}, we can state that $\Lambda^k(T^*M)$ is equipped with a smooth structure over $M$ via the canonical projection $\pi: \Lambda^k(T^*M) \to M$.

\noindent A smooth choice of one such alternating $k$-covector at every point of $M$ defines a differential $k$-form on the manifold.

\begin{definition}[Differential $k$-form]\label{def:k-forms}
	A \emph{differential $k$-form} on $M$ is a smooth section $\omega: M \to \Lambda^k(T^*M)$ of this space. We denote the $C^\infty(M)$-module of differential $k$-forms on $M$ by $\Omega^k(M)$.
\end{definition}

\begin{remark}\label{rem:forms-conventions}
	By convention, for $k = 0$ we set $\Lambda^0(T_p^*M) := \R$, which yields the identity $\Omega^0(M) = C^\infty(M)$. For $k = 1$, Definition \ref{def:k-forms} naturally recovers the space of $1$-forms $\Omega^1(M)$ introduced in Definition \ref{def:cot-space}. Furthermore, if $k > n = \dim M$, the alternating property implies that $\Lambda^k(T_p^*M) = \{0\}$, and consequently $\Omega^k(M) = \{0\}$.
\end{remark}


\subsection*{Operators on vector fields and forms}
With the structure of manifolds, tangent spaces, and the algebra of differential forms in place, we introduce the operators that act on them. In the geometric theory of mechanics, these operators map physical quantities between different spaces, through symmetries or coordinate transformations, for instance, and generate the dynamical equations. The pushforward maps velocity vectors. Its dual, the pullback, acts on configuration constraints and momenta.

First we see that there is a unique way to compute the product of vectors and covectors.

\begin{definition}[Interior product]
	Let $M$ be a smooth manifold, let $V \in \mathfrak{X}(M)$ be a smooth vector field, and let $\omega \in \Omega^k(M)$ be a smooth $k$-form. The \emph{interior product} of $\omega$ by $V$ is the $(k-1)$-form $\iota_V\omega \in \Omega^{k-1}(M)$ defined by
	\[
	(\iota_V\omega)(W_1,\dots,W_{k-1}) = \omega(V,W_1,\dots,W_{k-1}),
	\]
	for all vector fields $W_1,\dots,W_{k-1}$. For $k=0$, we set $\iota_V\omega = 0$.
\end{definition}

While the interior product acts as an algebraic contraction that reduces the degree of a form, we conversely require an operation to build higher-degree forms from lower-degree ones. To construct the full exterior algebra, we introduce an intrinsic product that combines a $k$-covector and an $l$-covector into a $(k+l)$-covector, rigorously preserving the alternating property.

\begin{definition}[Wedge product]\label{def:wedge-product}
	Let $M$ be a smooth manifold, and let $\omega \in \Lambda^k(T_p^*M)$ and $\eta \in \Lambda^l(T_p^*M)$ as defined in \ref{def:k-covector}. Their \emph{wedge product} (or exterior product) is the alternating $(k+l)$-covector $\omega \wedge \eta \in \Lambda^{k+l}(T_p^*M)$ defined by its evaluation on tangent vectors $v_1, \dots, v_{k+l} \in T_pM$:
	\[
	(\omega \wedge \eta)(v_1, \dots, v_{k+l}) := \frac{1}{k!l!} \sum_{\sigma \in S_{k+l}} \mathrm{sgn}(\sigma)\, \omega(v_{\sigma(1)}, \dots, v_{\sigma(k)}) \eta(v_{\sigma(k+1)}, \dots, v_{\sigma(k+l)}),
	\]
	where $S_{k+l}$ is the symmetric group of permutations of $\{1, \dots, k+l\}$ and $\mathrm{sgn}(\sigma)$ denotes the signature of the permutation. 
\end{definition}

\noindent This algebraic operation naturally extends pointwise to smooth sections, endowing the space of differential forms with the structure of a graded algebra.

\begin{remark}\label{rem:graded-commutativity}
	A fundamental algebraic consequence of Definition \ref{def:wedge-product} is its graded commutativity: for any $\alpha \in \Omega^k(M)$ and $\beta \in \Omega^l(M)$, the relation
	\[
	\alpha \wedge \beta = (-1)^{kl} \beta \wedge \alpha,
	\]
	holds identically. This entails that the wedge product of any $1$-form with itself strictly vanishes.
\end{remark}

We now examine the differential operators that map kinematic quantities between different manifolds. The most fundamental of such mappings actively transports tangent vectors along a smooth function.

\begin{definition}[Pushforward / Differential]\label{def:pushforward}
	Let $F: M \to N$ be a smooth map between smooth manifolds. For each $p \in M$, the \textit{pushforward} (or \textit{differential}) of $F$ at $p$ is the linear map 
	\[
		F_*: T_pM \to T_{F(p)}N,
	\]
	such that, for any tangent vector $X \in T_pM$ and for any smooth function $f \in C^\infty(N)$, is defined as
	\[
	(F_*X)(f) = X(f \circ F).
	\]
	The pushforward map is also denoted by $dF_p$.
\end{definition}

\begin{remark}
	We can state an alternative definition of \ref{def:pushforward} using the geometric description of tangent vectors as equivalence classes of curves. Let $\gamma \in C^\infty( (-\epsilon, \epsilon); M)$ such that $\gamma(0) = p$ and $\gamma'(0) = X$. Then the pushforward $F_*X$ is the velocity vector of the curve $F \circ \gamma$ at $t = 0$:
	\[
	F_*X = \frac{d}{dt}\Big|_{t=0} (F \circ \gamma)(t).
	\]
\end{remark}

The differential of a function yields the linear approximation of a scalar observable. To build the geometric machinery of classical mechanics, in particular the canonical symplectic form on the cotangent space, we need to extend this operation to higher-degree forms. The exterior derivative does exactly that.

\begin{definition}[Exterior derivative]
	Let $M$ be a smooth manifold, define a family of mappings $d^k(U) : \Omega^k(U) \to \Omega^{k+1}(U)$ with $k = 0,1,2,\dots,n$, and $U$ is open in $M$, which we denote by $d$, called the \emph{exterior derivative} on $M$, such that
	\begin{enumerate}
		\item[(i)] $d$ is a $\wedge$-antiderivation. That is, $d$ is $\mathbb{R}$ linear and for $\alpha \in \Omega^k(U)$, $\beta \in \Omega^l(U)$,
		\[
		d(\alpha\wedge\beta) = d\alpha\wedge\beta + (-1)^k\alpha\wedge d\beta,
		\]
		where $\wedge$ denotes the wedge-product of forms.
		\item[(ii)] If $f \in C^\infty(U)$, $d^0f = df$ is the differential as in Definition \ref{def:pushforward}.
		\item[(iii)] $d\circ d = 0$, that is, $d^{k+1}(U)\circ d^k(U) = 0$.
		\item[(iv)] $d$ is natural with respect to restrictions, that is, if $U \subset V \subset M$ are open and $\alpha \in \Omega^k(V)$, then $d(\alpha|_U) = (d\alpha)|_U$.
	\end{enumerate}
	It can be proved that such a family exist and it is unique.
\end{definition}

A crucial feature of the exterior derivative is its commutation with the pullback, a property that ensures the geometric consistency of coordinate transformations and symmetries. While the pushforward linearly transports velocities along a map, the pullback naturally acts in the opposite direction, transporting observables and differential forms.

\begin{definition}[Pullback]\label{def:pullback}
	Let $F: M \to N$ be a smooth map between smooth manifolds. For $k \geq 0$, let $\Omega^k(N)$ denote the space of smooth differential $k$-forms on $N$. The \textit{pullback} is the map
	\[
	F^*: \Omega^k(N) \to \Omega^k(M)
	\]
	defined as follows.
	\begin{enumerate}
		\item If $k = 0$, the space becomes $\Lambda = f \in C^\infty(N)$, the pullback $F^*f \in C^\infty(M)$ is defined by precomposition:
		\[
		F^*f = f \circ F.
		\]
		
		\item If $k \geq 1$, for $\Lambda \in \Omega^k(N)$, the pullback $F^*\Lambda \in \Omega^k(M)$ is defined pointwise, for each $p \in M$ and all tangent vectors $X_1, \dots, X_k \in T_pM$, by
		\[
		(F^*\Lambda)_p(X_1, \dots, X_k) = \Lambda_{F(p)}(F_*X_1, \dots, F_*X_k).
		\]
	\end{enumerate}
\end{definition}

\begin{remark}
	The case $k = 1$ recovers the pullback of covectors: for $\omega \in T_{F(p)}^*N$ and $X \in T_pM$,
	\[
	(F^*\omega)(X) = \omega(F_*X).
	\]
\end{remark}

Having described how differential forms and covectors transform under smooth maps, we now turn to the dynamical content of the vector field itself. In fact, a vector field generates, through its integral curves, a continuous family of diffeomorphisms of the manifold. This flow is the fundamental geometric object underlying the notion of time evolution, and it will later allow us to identify the trajectories of a Hamiltonian system with the flow of the corresponding Hamiltonian vector field.

\begin{definition}[Flow of a Vector Field]\label{def:flow}
	Let $X$ be a smooth vector field on a smooth manifold $M$. An \textit{integral curve} of $X$ is a smooth curve $\gamma: I \to M$, where $I \subseteq \R$ is an open interval, such that $\gamma'(t) = X(\gamma(t))$ for all $t \in I$.
	
	The \textit{(local) flow} of $X$ is a smooth map $\theta: \mathcal{D} \to M$, where $\mathcal{D} \subset \R \times M$ is an open neighborhood of $\{0\} \times M$ known as the \textit{flow domain}, satisfying the following properties:
	\begin{enumerate}
		\item For each $p \in M$, the slice $\mathcal{D}^{(p)} = \{t \in \R : (t, p) \in \mathcal{D}\}$ is an open interval containing $0$, and the curve $\theta^{(p)}: \mathcal{D}^{(p)} \to M$ defined by $\theta^{(p)}(t) = \theta(t, p)$ is the unique maximal integral curve of $X$ satisfying the initial condition $\theta(0, p) = p$, namely
		\[
		\frac{\partial \theta}{\partial t}(t, p) = X(\theta(t, p)).
		\]
		\item \textbf{One-Parameter Group Property:} Denoted by $\theta_t(p) = \theta(t, p)$, then for any $s \in \mathcal{D}^{(p)}$ and $t \in \mathcal{D}^{(\theta_s(p))}$ such that $t+s \in \mathcal{D}^{(p)}$, we have
		\[
		\theta_t(\theta_s(p)) = \theta_{t+s}(p).
		\]
	\end{enumerate}
	If $\mathcal{D} = \R \times M$, the vector field $X$ is called \textit{complete} and the maps $\theta_t: M \to M$ constitute a global one-parameter group of diffeomorphisms.
\end{definition}

The concept of flow translates the analytical problem of solving systems of differential equations into the geometric action of a one-parameter group of diffeomorphisms. In the context of dynamics, if $X$ is the Hamiltonian vector field governing a mechanical system, its integral curves represent the physical trajectories in phase space, and the associated flow $\theta_t$ corresponds to the time evolution operator, which inherently preserves the symplectic structure. This will be analysed later in \ref{sec:1-ham}

%INSERIRE DERIVATA DI LIE!!%


\section{Lie Algebras and Representations}
In this section we develop the algebraic component of classical mechanics. We first introduce Lie algebras in the abstract. We then specialize to Lie groups, smooth manifolds carrying a compatible group structure, and show how every Lie group determines a Lie algebra through its left-invariant vector fields, with the exponential map providing the bridge back from the algebra to the group. Finally, we introduce the notion of a representation, which realizes the abstract elements of a Lie algebra or Lie group as concrete linear operators on a vector space. This last step is what will allow us, in later chapters, to identify the continuous symmetries of a quantum system with observables that can actually be measured.

\textit{Note: In this chapter, $\K$ denotes an arbitrary commutative field with characteristic zero.}

\begin{definition}[Lie algebra]\label{def:lie-algebra}
	A \emph{Lie algebra} over $\mathbb{K}$ is a $\mathbb{K}$-vector space $\mathfrak{g}$ equipped with a bilinear map $[\cdot,\cdot]:\mathfrak{g}\times\mathfrak{g}\to\mathfrak{g}$ satisfying
	\begin{align}
	[X,Y] = -[Y,X], \qquad &\text{antisymmetry} \\
	[[X,Y],Z]+[[Y,Z],X]+[[Z,X],Y]=0, \qquad &\text{Jacobi identity}
	\end{align}
	for all $X,Y,Z\in\mathfrak{g}$. \cite{humphreysIntroductionLieAlgebras2010}
\end{definition}

\begin{remark}
	Every associative algebra $(\mathcal{A},\cdot)$ carries a canonical Lie algebra
	structure via the commutator $[a,b]:=ab-ba$. It satisfies identically the antisymmetry and the Jacobi identity:
	\[
		[[a,b],c] = (abc - bac) - c(ab - ba) = abc - bac +(- cab + cba) + (bca - bca) + (acb - acb) = -[[b,c],a] - [[c,a],b].
	\]
	The Jacobi identity is used in a Lie algebra to mimic the chain property of a derivation. For example, this allows $\mathfrak{g}$ to perfectly describe infinitesimal calculus on manifolds.
\end{remark}

\begin{example}[Vector fields as a Lie algebra]\label{ex:lie-fields}
	Let $M$ be a smooth manifold. The space of smooth vector fields $\mathfrak{X}(M)$, introduced in Definition \ref{def:vector-field}, forms an infinite-dimensional Lie algebra over $\R$. Since vector fields act as derivations on the algebra of smooth functions $C^\infty(M)$, we can define the Lie bracket of two vector fields $X, Y \in \mathfrak{X}(M)$ as their operator commutator:
	\[
	[X, Y](f) := X(Y(f)) - Y(X(f)), \qquad \forall f \in C^\infty(M).
	\]
	While the simple composition of two derivations $XY$ is not generally a derivation, their commutator inherently is, ensuring that $[X,Y] \in \mathfrak{X}(M)$ is again a well-defined vector field. The bilinearity and antisymmetry are manifest, and the Jacobi identity follows directly from the associative nature of operator composition.
\end{example}

\begin{example}
	One of the most important examples of Lie algebra is $\mathfrak{gl}(V)$, the space of all linear endomorphisms of a finite dimensional vector space $V$, equipped with the commutator $[A,B] = AB - BA$ as Lie bracket. When $V = \C^n$, $\mathfrak{gl}(V) \cong \mathfrak{gl}(n,\C)$ can be identified with $M(n,\C)$. A natural basis, denoted $E_{ij}$ for $i,j = 1,\dots,n$, is given by the \textit{matrix units}, \textit{i.e.}, the matrix with entry $1$ in position $(i,j)$ and $0$ elsewhere, satisfying $$[E_{ij}, E_{kl}] = \delta_{jk}E_{il} - \delta_{li}E_{kj}.$$
\end{example}

\begin{example}\label{ex:su-n}
	Another notable example of Lie algebra, as per Definition \ref{def:lie-algebra}, is $\mathfrak{su}(n)$. It consists of the matrices in $M(n, \C)$ which are traceless and anti-hermitian. The dimension of the resulting Lie algebra is $n^2 -1$. With $n=2$, the basis of $\mathfrak{su}(2)$ is $\{T_j = -i\sigma_j/2\}_{j=x,y,z}$, where $\sigma_j$ are the Pauli operators which are used for representing the spin of a fermion in quantum physics. Instead, with $n=3$, the basis is $\{-i\lambda_a/2\}_{a=1,...,8}$ which is composed by the eight Gell-Mann matrices $\lambda_a$ used in quantum chromodynamics.
\end{example}

We recall some basic concepts of the theory of Lie algebras.
A \textit{homomorphism} is a map between the same mathematical structures that preserves the inner operations. In particular, when the structure is a Lie algebra we can use the following definition.

\begin{definition}[Homomorphism of Lie Algebras]\label{def:lie-homo}
Let \(\mathfrak{g}\) and \(\mathfrak{h}\) be Lie algebras over the same field $\K$. A linear map \(\phi: \mathfrak{g} \rightarrow \mathfrak{h}\) is a Lie algebra homomorphism if it preserves the Lie bracket:
\[ 
	\phi([X, Y]) = [\phi(X), \phi(Y)] \quad \forall X, Y \in \mathfrak{g}.
\]
\end{definition}

\noindent When such a homomorphism is invertible, it identifies the two Lie algebras as abstractly the same structure.

\begin{definition}[Isomorphism of Lie Algebras]\label{def:lie-iso}
	An isomorphism is a Lie algebra homomorphism \(\phi: \mathfrak{g} \rightarrow \mathfrak{h}\) that is bijective.
\end{definition}

\subsection*{Lie Groups}
To benefit from the algebraic structure of Lie algebra in the framework of differential geometry, we must introduce manifolds equipped with compatible group operations. A tangent space exists at every point of a smooth manifold, but there is a restriction. In order to canonically endow the tangent space at a specific point with a Lie algebra structure, the manifold itself must be a Lie group.

To formalize this relationship, we first define a Lie group by requiring the group operations to be smooth.

\begin{definition}[Lie Group]\label{def:lie-group}
	A Lie group is a smooth manifold G together with a smooth map from $G \times G \to G$ that makes G into a group. Furthermore the inverse map $g \mapsto g^{-1}$ is smooth.
\end{definition}

\noindent We can specialize Definition \ref{def:lie-homo} in the case of Lie groups.

\begin{definition}[Homomorphism of Lie Groups]\label{def:lie-grhomo}
	A map $F: G \to H$ where G and H are Lie groups is called a homomorphism if it is a smooth map and a homomorphism between groups.
\end{definition}

Consider the tangent space of a Lie group. For every $g \in G$, we can define a map $L_g \in C^\infty(G;G)$ by $L_g(h) = gh$ called the \textit{"left multiplication"}. The differential $dL_g$ at a point $h \in G$ is a linear map from $T_h(G)$ to $T_{gh}(G)$ as stated in Definition \ref{def:pushforward}. A vector field X on G is called \textit{left-invariant} if
\[
	dL_g (X_h) = X_{gh} \qquad \forall g,h \in G.
\]
It is also true that, for any given vector $v \in T_e(G)$, there exists a unique smooth left-invariant vector field $X^v$. It can be constructed starting from the vector tangent to the identity by setting $X^v_e = v$. We can then use $L_g$ to define:
\[
	X^v_g = dL_g (X^v_e) = dL_g(v).
\]
By construction, this is left-invariant and smooth.
It is not difficult to show that the commutator of two left-invariant vector fields is, again, a left-invariant vector field. Also, it is worth noting that left-invariant vector fields form a linear subspace of $\mathfrak{M}$ closed respect to the commutator, called a \textit{Lie subalgebra} and thus they inherit the property of a Lie algebra. 

\begin{definition}[Lie algebra of a Lie group]\label{def:lie-gralg}
	The Lie algebra $\mathfrak{g}$ of a Lie group G is the tangent space at the identity with the bracket operation defined by
	\[
		[v, w] := [X^v, X^w]_e.
	\] 
\end{definition}

\noindent On a general smooth manifold, vector fields typically admit only local flows, potentially developing singularities in finite time. However, the underlying homogeneous structure of a Lie group circumvents this limitation. Local solutions to the associated differential equations can be indefinitely translated across the manifold via the group multiplication. This mechanism prevents finite-time divergences and guarantees the existence of global flows, ensuring that infinitesimal symmetries seamlessly integrate into one-parameter subgroups of finite transformations.

\begin{proposition}[Completeness of left-invariant vector fields]\label{prop:complete}
	Let $X^v \in \mathfrak{X}(G)$ be the left-invariant vector field associated to $v \in \mathfrak{g} = T_eG$. Then $X^v$ is complete: its maximal integral curve through $e$ is defined for all $t \in \R$.
\end{proposition}

Completeness alone does not yet reveal the full algebraic content of the flow. The following proposition shows that the integral curve of a left-invariant vector field is in fact a group homomorphism.

\begin{proposition}[One-parameter subgroup]\label{prop:one-param}
	Let $\gamma_v:\R\to G$ be the integral curve of $X^v$ through $e$. Then
	\[
	\gamma_v(t+s) = \gamma_v(t)\,\gamma_v(s) \qquad \forall\, t,s \in \R,
	\]
	$\gamma_v :\R \to G$ is a Lie group homomorphism called a \emph{one-parameter subgroup}.
\end{proposition}

These concepts naturally converge to establish a canonical correspondence between a Lie algebra and its associated Lie group. The mathematical tool enacting this exact correspondence is the exponential map.

\begin{definition}[Exponential map]\label{def:exp}
	The exponential map of $G$ is
	\[
	\exp: \mathfrak{g} \to G, \qquad \exp(v) := \gamma_v(1).
	\]
\end{definition}

\begin{proposition}[Properties of exponential map]\label{prop:exp-props}
	It can be shown that
	\begin{enumerate}
		\item $\exp(tv) = \gamma_v(t)$ $\forall t\in\R$.
		\item $\exp(0)=e$ and $d(\exp)_0 = \mathrm{id}_{\mathfrak g}$.
		\item Naturality: if $F:G\to H$ is a homomorphism of Lie groups with differential $dF_e:\mathfrak g\to\mathfrak h$, then
		\[
		F(\exp_G(v)) = \exp_H(dF_e(v)), \qquad v\in\mathfrak g.
		\]
	\end{enumerate}
\end{proposition}

The naturality property is particularly restrictive: it forces the exponential map to commute with any homomorphism of Lie groups. Before exploiting this, it is worth noting a subtlety in the global behaviour of $\exp$.

\begin{remark}
	Global surjectivity of $\exp$ is not automatic in general, it fails e.g.\ for $SL(2,\R)$, yet it is verified for $G$ compact and connected. Thus covers the classical groups used in this thesis ($SU(2)$, $U(n)$).
\end{remark}

\begin{remark}[Matrix Lie groups]
	For $G \subseteq GL(n,\C)$ a closed subgroup, $\mathfrak g \subset \mathfrak{gl}(n,\C) = M_n(\C)$ is identified with a linear subspace of matrices via $T_eG \subset T_eGL(n,\C) \cong M_n(\C)$. Left multiplication is matrix multiplication: $L_g(h) = gh$, so $dL_g(v) = gv$ for $v \in \mathfrak g$. The left-invariant vector field associated to $v$ is therefore $X^v_g = gv$, and its integral curve through $e$ solves the linear matrix ODE
	\[
	\dot\gamma(t) = \gamma(t)\, v, \qquad \gamma(0) = \Id.
	\]
	The unique solution is the matrix exponential series
	\[
	\gamma(t) = e^{tv} := \sum_{k=0}^\infty \frac{(tv)^k}{k!}.
	\]
	Hence $\exp(v) = e^{v}$: the abstract exponential map of Definition \ref{def:exp} coincides with the ordinary matrix exponential whenever $G$ is realized as a matrix group. In particular
	\[
	\mathfrak g = \{v \in M_n(\C) : e^{tv} \in G \ \forall t \in \R\}.
	\]
\end{remark}

\begin{example}[$\mathfrak{su}(2)$ from $SU(2)$ via the matrix exponential]
	Let $G = SU(2) = \{g \in GL(2,\C) : g^\dagger g = \Id,\ \det g = 1\}$, differentiating both defining conditions along $\gamma(t) = e^{tv}$, $\gamma(0)= \Id$, yields the linearized conditions on $v$:
	\[
	\frac{d}{dt}\Big|_{t=0} \big(e^{tv}\big)^\dagger e^{tv} = v^\dagger + v = 0, \qquad
	\frac{d}{dt}\Big|_{t=0} \det e^{tv} = \operatorname{tr}(v) = 0,
	\]
	using $\det e^{tv} = e^{t\,\operatorname{tr} v}$. It follows that
	\[
	\mathfrak{su}(2) = \{v \in M_2(\C) : v^\dagger = -v,\ \operatorname{tr}(v) = 0\},
	\]
	a $3$-dimensional real vector space. Take $v = -\tfrac{i\theta}{2}\sigma_3$, $\theta \in \R$, $\sigma_3 = \mathrm{diag}(1,-1)$. Since $\sigma_3^2 = \Id$, the series splits into even and odd powers:
	\[
	\exp(v) = e^{-i\theta\sigma_3/2} = \cos\!\left(\frac{\theta}{2}\right) \Id - i\sin\!\left(\frac{\theta}{2}\right)\sigma_3
	= \begin{pmatrix} e^{-i\theta/2} & 0 \\ 0 & e^{i\theta/2} \end{pmatrix}.
	\]
	This matrix is unitary with determinant $1$ for every $\theta$, confirming $\exp(v) \in SU(2)$ and $\theta \mapsto \exp(v)$ satisfies the one-parameter subgroup law, consistently with Proposition \ref{prop:one-param}. 
\end{example}


\subsection*{Representations of groups}
Lie groups and Lie algebras are abstract mathematical structures. To use them in physics, we need to identify their elements with linear operators acting on a vector space. This identification is called a representation. In quantum mechanics and quantum information theory the relevant state space is a Hilbert space, and representation theory gives us the dictionary between abstract symmetry generators and concrete linear operators. This is what turns continuous symmetries into observables we can actually measure. 

In this section, we start with the general definition of a representation for arbitrary groups, then specialize to Lie groups and Lie algebras. \cite{hallLieGroupsLie2003}

\begin{definition}[Representation of a group]\label{def:group-rep}
	Let G be a group and V a vector space over $\K$. A representation of G on V is a homomorphism between groups:
	\[
		\rho : G \to GL(V),
	\]
	where $GL(V)$ denotes the general linear group, the set of all the automorphisms $\psi: V \to V$. A representation is often denoted with $(V, \rho)$.
\end{definition}


\noindent In the case of Lie algebras this can be transported as the following.

\begin{definition}[Representation of a Lie Algebra]\label{def:lie-rep}
	A representation of a Lie Algebra $\mathfrak{g}$ is a homomorphism $\phi: \mathfrak{g} \to \mathfrak{gl}(V)$ where V is a $\K$-vector space.
\end{definition}

\begin{remark}
	It is worth noting that this definition is slightly different from the previous one. In the case of a group, the range of the homomorphism $\rho$ is the set of all automorphisms. Instead in Definition \autoref{def:lie-rep}, the maps onto $\mathfrak{gl}(V)$ are not all isomorphisms. Observe that $\mathfrak{gl}(V)$ is the \emph{tangent space} to $GL(V)$ at the identity.  Defining an exponential map $g(t) = e^{tX}$ one has $g'(0) = X$, which is generically non-invertible even though every $g(t)$ is invertible. Algebraically, this reflects the fact that a Lie algebra homomorphism $\rho : \mathfrak{g} \to \mathfrak{gl}(V)$ need only preserve the bracket, $\rho([X,Y]) = [\rho(X),\rho(Y)]$, whereas a group homomorphism $\rho : G \to GL(V)$ must preserve inverses, which forces $\rho(g) \in GL(V)$ for every $g \in G$.
\end{remark}

\begin{example}
	We define the \textit{adjoint representation} as the homomorphism $\ad: \mathfrak{g} \to \mathfrak{gl}(\mathfrak{g})$ which sends $x \mapsto \ad{x}$ with $\ad{x}$ being an endomorphism on $\mathfrak{g}$ defined as \[ \ad{x}(y) := [x,y].\]
	In addition, it can be easily proved that \[
		\ad{[x,y]} = [\ad{x},\ad{y}],
	\]
	showing the clear preservation of Lie brackets.
\end{example}

One of the most relevant properties of a representation is faithfulness.
\begin{definition}[Faithful Representation]\label{def:rep-faith}
	A representation $(V, \rho)$ is said to be \textit{faithful} if $\rho$ is injective that is
	\[
		\ker{\rho} = \{g \in G \ | \ \rho(g) = \Id_V\} = \{e\},
	\]
	where $e$ is the identity element of the group.
\end{definition}

Establishing whether a representation is faithful determines if it captures the full algebraic structure of the original group without redundancy. However, physical systems may realize the identical abstract symmetry through fundamentally different matrix formulations, reflecting a diverse choice of basis for the underlying state space. Therefore, before attempting to systematically decompose these spaces, it is strictly necessary to formulate a formal criterion to identify when two representations are mathematically equivalent and physically indistinguishable.

\begin{definition}[Equivalent representations] \label{def:rep-equiv}
	Let $(V, \rho)$ and $(W, \pi)$ be two representations of a group $G$, with V and W two $\K$-vector spaces. The representations are said to be equivalent if there exists a vector space isomorphism $T: V \to W$ such that:
	\[
		T\rho(g) = \pi(g)T \qquad \forall g \in G.
	\]
\end{definition}

\begin{remark}
	Another way to express Definition \autoref{def:rep-equiv} is the following. In order to define a representation equivalent to $(V, \rho)$ on a $\K$-vector space W, given a vector space isomorphism $T$ we define the homomorphism of the equivalent representation as
	\[
	\pi := T \circ \rho \circ T^{-1}.
	\]
\end{remark}

\noindent Now we are ready to define what it means to decompose representations into the so called \textit{irreducible representations}.

\begin{definition}[Irreducibility of a representation] \label{def:rep-irreps}
	Let (V, $\pi$) be a finite-dimensional representation of a group G, acting on a space V. 
	We say that:
	\begin{itemize}
	 \item A subspace W of V is invariant if $\pi(g)(w) \in W \quad \forall w \in W , \forall g \in G$. 
	 \item An invariant subspace W is nontrivial if $W \neq \{0\}$ and $W \neq V$. 
	 \item A representation with no nontrivial invariant subspaces is irreducible.
	\end{itemize}
	The terms invariant, nontrivial, and irreducible are defined analogously for representations of Lie algebras.
\end{definition}

\noindent Thanks to such a definition we can state a notable result.

\begin{lemma}[Schur's Lemma]\label{lem:schur}
	Let $\pi: G \to GL(V)$ be a finite-dimensional \emph{irreducible representation}  over $\C$.
	\begin{enumerate}
		\item If $T: V \to V$ is a linear map commuting with $\pi(g)$ for every $g$, then $T = \lambda \, \mathrm{id}_V$ for some $\lambda \in \C$.
		\item More generally, if $\pi_1: G\to GL(V_1)$, $\pi_2: G\to GL(V_2)$ are both irreducible and $T: V_1 \to V_2$ satisfies $T\pi_1(g) = \pi_2(g)T$ for all $g$, then either $T=0$ or $T$ is an isomorphism.
	\end{enumerate}
\end{lemma}

The profound physical implication of Schur's Lemma manifests prominently in the study of Casimir operators. Those are generically non-linear algebraic invariants constructed from the generators of the Lie algebra that, by definition, commute with all elements of the algebra itself. Since a Casimir operator commutes with the entire representation, Schur's Lemma dictates that it must act as a scalar multiple of the identity on any irreducible representation space. This scalar value ultimately furnishes the exact quantum numbers necessary to classify the fundamental multiplets of a quantum system. This lemma can be extended to Lie Algebras as well, we can show an example of its application on $\mathfrak{su}(2)$.

\begin{example}[Schur's Lemma on $\mathfrak{su}(2)$]
	Let $\pi: \mathfrak{su}(2) \to \mathrm{End}(V)$ be a finite-dimensional irreducible representation, let $J_i := i T_i$ for $i = x,y,z$ as shown in Example \ref{ex:su-n} and define:
	\[
	\hat J^2 := \hat J_x^2 + \hat J_y^2 + \hat J_z^2
	\]
	be the total angular momentum squared operator. Using $[\hat J_i, \hat J_j] = i\varepsilon_{ijk}\hat J_k$, a direct computation yields $[\hat J^2, \hat J_i] = 0$ for every $i$. Hence $\hat J^2$ commutes with the entire representation.
	
	\noindent By Schur's Lemma \ref{lem:schur} item 1., any operator commuting with an irreducible representation is a scalar multiple of the identity:
	\[
	\hat J^2 = \lambda\, \Id, \qquad \text{for some } \lambda \in \R.
	\]
	
	\noindent In this simple case we see how Schur's Lemma guarantees $\hat J^2$ takes one fixed value on the whole representation, not just on individual states.
\end{example}

\noindent To conclude, we state an important definition that will be needed in the next chapters.

\begin{definition}[Strongly continuous unitary representation] \label{def:rep-scur}
	Let $\H$ be a Hilbert space, and let $U(\H)$ denote the group of unitary operators on $\H$. Then, a homomorphism $\Pi \colon G \to U(\mathcal{H})$, where G is a Lie group, is called a strongly continuous unitary representation of $G$ if $\Pi$ satisfies the following continuity condition called strong continuity: 
	\[
 		A_n, A \in G \quad \text{and} \quad  A_n \to A \quad \implies \quad \Pi(A_n) (v) \to \Pi(A) (v) \quad \forall v \in \H.
 	\] 
\end{definition}

\section{Elements of symplectic geometry}

In the previous sections, we established the formalism of differential geometry and the theory of Lie groups, which are essential tools for describing the global symmetries and the topology of physical systems. However, the manifold describing just the positions of the particles, called \textit{the configuration manifold} turns out to be insufficient. Dynamical evolution naturally arises from the phase space, mathematically identified with the cotangent space of the base manifold. This structure is complete enough to yield a way to state an isomorphism between covectors and vector fields. The framework provided is called symplectic geometry. 

In this section, we will introduce the foundations of symplectic manifolds, starting from the fundamental notion of a closed and non-degenerate $2$-form, which constitutes the mathematical core required to define canonical transformations and govern Hamiltonian dynamics.

\begin{definition}[Non-degenerate 2-form on a manifold]\label{def:non-deg}
	Let $M$ be a smooth manifold and $\omega \in \Omega^2(M)$. We say $\omega$ is \emph{non-degenerate} if, for every $p \in M$, the bilinear form $\omega_p$ is non-degenerate on $T_pM$: for every $X \in T_pM$,
	\[
	\big(\omega_p(X,Y) = 0 \ \ \forall Y \in T_pM\big) \implies X = 0.
	\]
\end{definition}

The next proposition shows that any skew-symmetric two-form, degenerate or not, admits a pointwise normal form, a purely algebraic fact from which the non-degenerate case will follow as an immediate specialization.

\begin{proposition}\label{prop:skew-normal-form-manifold}
	Let $\omega \in \Omega^2(M)$ on a smooth manifold $M$, and for $p \in M$ set $r_p := \mathrm{rank}(\omega_p)$. Then for every $p \in M$ there exists an integer $n_p$ with $r_p = 2n_p$, and an ordered basis $\hat e = (e_i)$ of $T_pM$, with dual ordered basis $(\alpha^i)$, such that
	\[
	\omega_p = \sum_{i=1}^{n_p} \alpha^i \wedge \alpha^{i+n_p},
	\]
	or, equivalently, the matrix of $\omega_p$ in this basis is
	\[
	\begin{pmatrix}
		0 & \Id & 0 \\
		-\Id & 0 & 0 \\
		0 & 0 & 0
	\end{pmatrix},
	\]
	where $\Id$ is the $n_p \times n_p$ identity matrix and the bottom-right block has size $(d-2n_p)\times(d-2n_p)$, $d := \dim T_pM$.
\end{proposition}

\begin{remark}
	This is pointwise linear algebra: $(e_i)$ is a basis of $T_pM$ alone, chosen independently at each $p$, with no claim that $p \mapsto e_i(p)$ varies smoothly, or even continuously.
\end{remark}

\noindent Specializing the previous proposition to the non-degenerate case removes the degenerate block entirely and fixes the rank at every point, with the following immediate consequences.

\begin{corollary}\label{cor:omega-prop}
	Suppose $\omega$ is non-degenerate by Definition \ref{def:non-deg} , then $r_p = \dim T_pM =: d$ for every $p \in M$, so the bottom-right block above has no size, and $n_p = d/2$ is the \emph{same} integer at every point, call that $n$. In particular:
	\begin{itemize}
		\item $d$ must be even: a skew-symmetric non-degenerate $2$-form can exist only on an even-dimensional manifold (on an odd-dimensional one, every $2$-form is degenerate at every point).
		\item we can define $\omega^n := \omega \wedge \cdots \wedge \omega$ with $n=d/2$ factors, which is nowhere-vanishing and, in the basis used above, it can be written as $\omega_p^{\,n} = n!\,\alpha^1\wedge\alpha^{n+1}\wedge\cdots\wedge\alpha^n\wedge\alpha^{2n} \neq 0$. So, $\omega^n$ is a volume form on $M$, that can be standardized as
		\[
		\Omega_\omega := \frac{(-1)^{n(n-1)/2}}{n!} \omega^n.
		\]
	\end{itemize}
\end{corollary}

Beyond restricting the dimension and providing a canonical volume measure, the non-degeneracy of $\omega$ plays a far more active role in the formulation of mechanics. Specifically, a non-degenerate bilinear form naturally induces a point-wise vector space isomorphism between the tangent and cotangent spaces. This duality is the fundamental mathematical machinery required to translate the differential of an observable (such as the Hamiltonian) into a unique vector field that generates the dynamics.

\begin{definition}[Musical isomorphisms]\label{def:musical}
	Let M be a smooth manifold and $\omega \in \Omega^2(M)$ be non-degenerate by Definition \ref{def:non-deg}. Then define the map $\flat :\mathfrak{X}(M) \to \Omega^1(M)$ so that a vector field $$X \mapsto X^\flat = \iota_X\omega = \omega(X, \cdot).$$ 
	
	Let instead the inverse map be $\sharp: \Omega^1(M) \to \mathfrak{X}(M)$ so that the form $\alpha \mapsto \alpha^\sharp$ where $\alpha^\sharp$ is the only vector field $X$ such that $\iota_X \omega= \alpha$.
\end{definition}

\begin{remark}
	Note that the inverse map $\sharp$ is well-defined because $\flat$ is bijective being $\omega$ non-degenerate as in Definition \ref{def:non-deg}. We can clearly see that $(X^\flat)^\sharp = X$ and $(\alpha^\sharp)^\flat = \alpha$ due to the non-degeneracy of $\omega$.
\end{remark}

\noindent With the musical isomorphisms in hand, we are ready to state a central theorem that underlines the importance of the closure of the symplectic form.

\begin{theorem}[Darboux Theorem]\label{thm:darboux}
	Suppose $\omega$ is a non-degenerate two-form on a 2n-manifold M by Definition \ref{def:non-deg}. Then $\omega$ is closed, meaning $d \omega = 0$ (where $d$ denotes the exterior derivative), if and only if there is a chart $(U, \phi)$ at each $m \in M$ such that $\phi(m) = \boldsymbol{0}$ and $\phi(u) = (x^1(u),..,x^n(u), y^1(u),...,y^n(u))$ having
	\[
		\omega|_U = \sum_{i=1}^n dx^i \wedge dy^i.
	\]
\end{theorem}

\begin{remark}\label{rem:darboux-consequences}
	One direction of the equivalence is easy. If we start from the fact that the canonical local expression of the form is exact, then it is automatically closed. The true analytical depth of Darboux's Theorem, however, lies in its sufficiency. It proves that the closure of $\omega$ is the only requirement needed to guarantee the existence of a local canonical frame. Consequently, any two symplectic manifolds of the same dimension are locally diffeomorphic.
\end{remark}

\noindent Finally, we can state the next important definition that encodes all the structure to derive Hamiltonian mechanics in the natural way.

\begin{definition}[Symplectic manifold]\label{def:sympl-man}
	A symplectic form $\omega$ on a smooth manifold M is a non-degenerate, closed, two-form on M. A symplectic manifold $(M, \omega)$ is a smooth manifold M together with a symplectic form $\omega$ on M. The charts guaranteed by the Darboux theorem \ref{thm:darboux} are called symplectic charts and the coordinate are called \emph{canonical coordinates}.
\end{definition}

Having in hand the geometric framework of the phase space, it is natural to investigate what class of transformations preserves the structure of the symplectic manifold. Such mappings are fundamental as they generate all the symmetries within Hamiltonian mechanics. 

\begin{definition}[Canonical transformation]\label{def:canonical-transformation}
	Let $(M,\omega)$ and $(N,\rho)$ be symplectic manifolds, F a $C^\infty$-mapping $F:M \to N$ is called \emph{symplectic} or \emph{canonical transformation} if $F^*\rho = \omega$.
\end{definition}

\noindent Since canonical transformations carry $\rho$ exactly onto $\omega$, they must preserve the associated volume form, which forces them to be local diffeomorphisms.


\begin{proposition}
	If $(M,\omega)$ and $(N,\rho)$ be symplectic manifolds with $\dim M = \dim N$ and $F:M \to N$ symplectic, then F is volume preserving, $$F^*\Omega_\rho=\Omega_\omega,$$ with $\Omega_\omega$ defined as in \ref{cor:omega-prop}, and F is a local diffeomorphism.
\end{proposition}

The next theorem is notationally dense, but it is central to this introduction to symplectic geometry. The problem is the following: Constructing a 1-form on a manifold M, requires a rule that assign a real number to any tangent vector. When M is already a cotangent space of another manifold Q (as in the case of the phase space in Hamiltonian mechanics), there is a natural way to accomplish this.

\medskip 
Fix a $q \in Q$, from the perspective of Q, every point in M is a 1-form, which we denote as $\alpha_q$. From the perspective of M, we can instead consider the tangent vector at the point $\alpha_q$, calling it $w_{\alpha_q}$. 
At this point, the natural way to create a 1-form is to use $\alpha_q$ viewed as a form acting on tangent vectors in $TQ$. Thus, we can apply a pushforward to $w_{\alpha_q}$ in order to move it to the tangent space of Q and finally pairing it with $\alpha_q$. This is how the \textit{canonical form} $\theta_0$ works.

\begin{theorem}[Canonical forms]\label{thm:canon}
	Let Q be an $n$-manifold and $M = T^*Q$. Consider now the maps $\tau^*_Q: T^*Q \to Q$ and $d\tau^*_Q: TM \to TQ$, which is the pushforward of $\tau^*_Q$. Let $\alpha_q$ denote a point in $T^*Q$, with $q \in Q$, and $w_{\alpha_q} \in TM$ in the fiber above $\alpha_q$. Define
	\[
		\theta_{\alpha_q} : T_{\alpha_q}M \to \R \qquad w_{\alpha_q} \mapsto \alpha_q(d \tau^*_Q(w_{\alpha_q})),
	\] 
	and the map $\theta_0: \alpha_q \to \theta_{\alpha_q}$. In this case $\theta_0 \in \Omega^1(M)$ and $\omega_0=-d\theta_0$ is a symplectic form on M. The forms $\theta_0$ and $\omega_0$ are called the canonical forms on M.
\end{theorem}

\begin{remark}
	The superscript $*$ in $\tau^*_Q$ there simply labels it as the \emph{cotangent}-space projection, as opposed to $\tau_Q: TQ \to Q$. It is not the pullback operator. 
\end{remark}

\begin{corollary}[Canonical forms in coordinates]\label{crl:canon-coord}
	From Theorem \ref{thm:canon}, in finite-dimensional smooth manifolds written in local coordinates of $M$ $(q^1,...,q^n,p_1,...,p_n)$ , it follows that
	\[
		\theta_0 = \sum_{i=1}^n p_i dq^i \qquad \omega_0 = \sum_{i=1}^n dq^i \wedge dp_i.
	\]
\end{corollary}

\noindent We can also prove that these canonical forms are not changed by any diffeomorphism on the manifold.
\begin{theorem}[Invariance of the canonical forms]\label{thm:inv-canon}
	Let Q be a smooth manifold and $f: Q \to Q$ a diffeomorphism. Define the lift of f as
	\[
		T^*f : T^*Q \to T^*Q, \qquad (T^*f(\alpha_q))(v) = \alpha_q(df(v)) \quad \forall v \in T_qQ.
	\]
	Then $T^*f$ is symplectic and the canonical form $\theta_0$, defined in Theorem \ref{thm:canon}, is invariant under the application of the pulled-back lift:  $(T^*f)^*\theta_0 = \theta_0$.
\end{theorem}

\begin{remark}
	%The lift of diffeomorphisms are defined to make the following diagram commute.
	%\[
	%\begin{tikzcd}
	%	T^*Q \arrow[r, "T^*f"] \arrow[d, "\tau_Q^*"'] & T^*Q \arrow[d, "\tau_Q^*"] \\
	%	Q & Q \arrow[l, "f"]
	%\end{tikzcd}
	%\]
	In coordinates, if $f(q^1,...q^n) = (Q^1,...,Q^n)$, then the lift of $f$ is a map that 
	\[
		(Q^1,...,Q^n, P_1,...,P_n) \mapsto (q^1,...,q^n, p_1,...,p_n),
	\] 
	where
	\[
		p_j = \sum_{i=1}^n \frac{\partial Q^i}{\partial q^j} P_i.
	\]
	%Notice also that
	%\[
	%T^*(f \circ g) = T^*g \circ T^*f
	%\]
	%that is a law of composition similar to $T(f \circ g) = Tf \circ Tg$.
\end{remark}


\printbibliography

\end{document}
